\section{Experiments}
\label{sec:experiments}

\subsection{Setup}

Detection uses \DetectorName{} with CLIP embeddings per instance; text
conditions use \LLMName{} and the image-conditioned baseline uses \VLMName,
both served through vLLM. Change detection is evaluated on LEVIR-CD (test, 128
tiles) and WHU-CD (test, 690 pairs); report quality on LEVIR-CC. All timings
are measured on one A100 80GB.

\textbf{Integrity.} Every run asserts and records four properties: train and
test pair identifiers are disjoint; their parent scenes are disjoint, so two
crops of one tile cannot straddle the split; all thresholds are selected on the
validation split, never on test; and prompt style examples, where used, are
drawn only from the training split. The assertions are part of the result
files, not of the prose.

\textbf{Tiers.} Methods trained on a benchmark's own training split and methods
that have never seen it are not competing on equal terms, and printing them in
one block is misleading whichever way the numbers fall. Tables group by tier
and label each group. Our LEVIR-CD result occupies a third position we state
rather than blur: the detector never saw LEVIR-CD, but the pairing head was
trained on its training split from mask-derived labels, so it is not
zero-shot. Section~\ref{ssec:e2} is the zero-shot claim.

\subsection{Change Detection (E1)}

\begin{table}[t]
\centering
\caption{Pixel-level change detection on LEVIR-CD. Baselines are fully supervised; ours is open-vocabulary with a 20k-parameter pairing head.}
\label{tab:levir_cd_pixel}
\setlength{\tabcolsep}{4pt}
\begin{tabular}{lcccc}
\toprule
Method & P & R & F1 & IoU \\
\midrule
\multicolumn{5}{l}{\emph{Trained on this dataset (reference ceiling)}} \\
FC-Siam-Diff~\cite{daudt2018fully} & 89.53 & 83.31 & 86.31 & 75.92 \\
STANet~\cite{chen2020spatial} & 83.81 & 91.00 & 87.26 & 77.40 \\
SNUNet~\cite{fang2021snunet} & 89.18 & 87.17 & 88.16 & 78.83 \\
BIT~\cite{chen2021remote} & 89.24 & 89.37 & 89.31 & 80.68 \\
ChangeFormer~\cite{bandara2022transformer} & 92.05 & 88.80 & 90.40 & 82.48 \\
TinyCD~\cite{codegoni2023tinycd} & -- & -- & 91.05 & 83.57 \\
\midrule
\multicolumn{5}{l}{\emph{Zero-shot / open-vocabulary}} \\
AnyChange (SAM zero-shot)~\cite{zheng2024segment} & 13.70 & 83.00 & 23.40 & -- \\
AnyChange-H + 3-point query~\cite{zheng2024segment} & 28.50 & 81.10 & 42.20 & -- \\
\midrule
\multicolumn{5}{l}{\emph{This work}} \\
geo-only & 14.34 & 64.73 & 23.47 & 13.30 \\
heuristic (production) & 16.08 & 62.67 & 25.59 & 14.67 \\
learned head + greedy & 56.35 & 51.84 & 54.00 & 36.99 \\
learned head, no verifier & 12.49 & 76.66 & 21.48 & 12.03 \\
learned head, no cross-frame & 47.26 & 51.52 & 49.30 & 32.71 \\
\textbf{learned head (ours)} & \textbf{56.22} & \textbf{51.97} & \textbf{54.01} & \textbf{36.99} \\
\bottomrule
\end{tabular}
\end{table}

\begin{table}[t]
\centering
\caption{Instance-level change detection on LEVIR-CD (IoU $\geq$ 0.5).}
\label{tab:levir_cd_instance}
\setlength{\tabcolsep}{4pt}
\begin{tabular}{lccc}
\toprule
Method & P & R & F1 \\
\midrule
\multicolumn{4}{l}{\emph{This work}} \\
geo-only & 32.36 & 39.60 & 35.62 \\
heuristic (production) & 34.11 & 37.48 & 35.72 \\
learned head + greedy & 80.86 & 40.85 & 54.28 \\
learned head, no verifier & 24.34 & 43.62 & 31.25 \\
learned head, no cross-frame & 64.97 & 41.38 & 50.56 \\
\textbf{learned head (ours)} & \textbf{80.79} & \textbf{40.87} & \textbf{54.28} \\
\bottomrule
\end{tabular}
\end{table}


The row worth pointing at is not ours. Geometry-only pairing (\EOneGeoFOne),
the production heuristic (\EOneHeuristicFOne) and AnyChange
(\AnyChangeFOne) all land within two points of each other, at precision 13--16
with high recall. Methods that do not learn to pair converge on the same place
regardless of segmentation quality, and the \HeadParamsRound-parameter head is
what leaves that cluster: \EOnePixelFOne{} pixel F1, $2.3\times$ AnyChange's,
above even its human-assisted three-point setting (\AnyChangeHFOne).

Both ablations are load-bearing. Without the verifier the system scores
\EOneNoVerifierFOne{} --- below geometry-only --- because the head proposes
matches and something has to reject them. Retrained without cross-frame
evidence it loses \EOneCrossFrameGain{} pixel F1.

\subsection{Transfer Without Retraining (E2)}
\label{ssec:e2}

\begin{table}[t]
\centering
\caption{Pixel-level change detection on WHU-CD. Baselines are fully supervised; ours is open-vocabulary with a 20k-parameter pairing head.}
\label{tab:whu_cd_pixel}
\setlength{\tabcolsep}{4pt}
\begin{tabular}{lcccc}
\toprule
Method & P & R & F1 & IoU \\
\midrule
%% TODO fill and verify in baselines.json: FC-Siam-Diff [supervised], BIT [supervised], ChangeFormer [supervised], AnyChange (SAM zero-shot) [zero-shot]
\midrule
\multicolumn{5}{l}{\emph{This work}} \\
geo-only & 23.35 & 42.57 & 30.16 & 17.76 \\
heuristic (production) & 21.73 & 42.98 & 28.86 & 16.87 \\
learned head + greedy & 73.06 & 48.05 & 57.97 & 40.82 \\
learned head, no verifier & 12.13 & 53.50 & 19.78 & 10.98 \\
\textbf{learned head (ours)} & \textbf{72.89} & \textbf{48.05} & \textbf{57.92} & \textbf{40.77} \\
\bottomrule
\end{tabular}
\end{table}

\begin{table}[t]
\centering
\caption{Instance-level change detection on WHU-CD (IoU $\geq$ 0.5).}
\label{tab:whu_cd_instance}
\setlength{\tabcolsep}{4pt}
\begin{tabular}{lccc}
\toprule
Method & P & R & F1 \\
\midrule
\multicolumn{4}{l}{\emph{This work}} \\
geo-only & 4.73 & 55.62 & 8.73 \\
heuristic (production) & 4.93 & 55.93 & 9.06 \\
learned head + greedy & 44.18 & 63.83 & 52.22 \\
learned head, no verifier & 2.51 & 68.90 & 4.85 \\
learned head, no cross-frame & 38.63 & 63.32 & 47.98 \\
\textbf{learned head (ours)} & \textbf{44.18} & \textbf{63.83} & \textbf{52.22} \\
\bottomrule
\end{tabular}
\end{table}


The head is applied to WHU-CD unchanged: no retraining, and the thresholds are
those selected on LEVIR-CD's validation split. It reaches \ETwoPixelFOne{}
pixel F1 against \ETwoHeuristicFOne{} for the heuristic it replaces, and
\ETwoInstFOne{} instance F1 against \ETwoHeuristicInstFOne. The verifier
ablation collapses below geometry-only exactly as in E1.

The instance columns state it more usefully than F1 does. Against \ETwoNGT{}
ground-truth instances the heuristic emits \ETwoHeuristicNPred{} predictions
and ours emits \ETwoNPred: count error falls from \ETwoHeuristicCountMAE{} to
\ETwoCountMAE. For a service that reports how many buildings appeared, that is
the difference between a usable number and noise.

\textbf{This is not evidence that transfer is easier than the source domain.}
WHU-CD test averages \ETwoGTDensity{} ground-truth instances per tile against
LEVIR-CD's \EOneGTDensity, a $38\times$ difference in change density, so
\ETwoPixelFOne{} and \EOnePixelFOne{} are not comparable as difficulties. We
report both densities for that reason.

\subsection{Grounding and Factuality (E3, E4)}

\begin{table}[t]
\centering
\caption{Report factuality by grounding condition (caption style). CFS is claim-level precision/recall against the human references; Hal is the share of generated claims with no ground-truth support.}
\label{tab:factuality_caption}
\setlength{\tabcolsep}{4pt}
\begin{tabular}{lccccc}
\toprule
Method & CFS-P & CFS-R & CFS-F1 & Hal & ChgAcc \\
\midrule
\multicolumn{6}{l}{\emph{This work}} \\
template & 61.55 & 39.60 & 48.20 & 38.45 & 76.46 \\
vlm_direct & 34.97 & 24.03 & 28.49 & 65.03 & 45.78 \\
llm_raw & 55.87 & 38.06 & 45.28 & 44.13 & 70.40 \\
llm_struct & 61.73 & 39.76 & 48.37 & 38.27 & 76.46 \\
llm_flat_rag & 61.96 & 39.85 & 48.51 & 38.04 & 76.57 \\
\textbf{llm_graphrag} & \textbf{61.72} & \textbf{39.70} & \textbf{48.32} & \textbf{38.28} & \textbf{76.52} \\
\bottomrule
\end{tabular}
\end{table}

\begin{table}[t]
\centering
\caption{Change captioning on LEVIR-CC. Baselines are trained on the LEVIR-CC training split; our conditions are zero-shot apart from the pairing head.}
\label{tab:levir_cc_caption}
\setlength{\tabcolsep}{3pt}
\begin{tabular}{lccccc}
\toprule
Method & BLEU-1 & BLEU-4 & METEOR & ROUGE-L & CIDEr-D \\
\midrule
\multicolumn{6}{l}{\emph{Trained on this dataset (reference ceiling)}} \\
DUDA~\cite{park2019robust} & 81.44 & 57.79 & 37.15 & 71.04 & 124.32 \\
MCCFormers-S~\cite{qiu2021describing} & 82.16 & 59.41 & 38.26 & 72.10 & 128.34 \\
MCCFormers-D~\cite{qiu2021describing} & 80.49 & 57.34 & 38.23 & 71.40 & 126.85 \\
RSICCformer~\cite{liu2022remote} & 84.11 & 61.93 & 38.79 & 73.02 & 131.40 \\
Chg2Cap~\cite{chang2024changes} & 86.14 & 64.39 & 40.03 & 75.12 & 136.61 \\
\midrule
\multicolumn{6}{l}{\emph{This work}} \\
\texttt{template} & 50.33 & 26.03 & 51.50 & 47.10 & 25.29 \\
\texttt{vlm\_direct} & 31.85 & 4.51 & 21.71 & 25.27 & 11.47 \\
\texttt{llm\_raw} & 53.02 & 29.15 & 40.66 & 44.72 & 60.11 \\
\texttt{llm\_struct} & 48.09 & 28.15 & 49.90 & 51.49 & 77.77 \\
\texttt{llm\_flat\_rag} & 49.52 & 29.19 & 50.42 & 52.06 & 78.35 \\
\textbf{\texttt{llm\_graphrag}} & \textbf{49.53} & \textbf{29.56} & \textbf{49.81} & \textbf{51.63} & \textbf{78.00} \\
\bottomrule
\end{tabular}
\end{table}


Two results in this table have to be reported rather than buried.

\textbf{The external baseline decides at chance.} \texttt{vlm\_direct} reaches
ChgAcc \EFourVLMChgAcc\% --- it is guessing whether anything changed --- while
producing confident, well-formed prose. CFS-F1 \EFourVLMFOne{} against
\EFourGraphRAGFOne{} for the grounded pipeline.

\textbf{Aggregation is worth \EFourAggGain{} points; retrieval and graph
aggregation on top of it are worth nothing measurable.} Structuring the raw
detection dump into a change inventory moves CFS-F1 from
\EFourRawFOne{} (\texttt{llm\_raw}) to \EFourStructFOne{}
(\texttt{llm\_struct}). Adding top-$k$ retrieval of past observations
(\EFourFlatRagFOne) or replacing it with graph entity and community
aggregation (\EFourGraphRAGFOne) moves it by less than
\EFourLadderSpread{} points, and not consistently in the graph's favour:
over \EFourNCrops{} crops \texttt{llm\_flat\_rag} finishes marginally
ahead. We read the three as indistinguishable rather than ranked.

The reason is structural rather than a shortcoming of the graph. All three
conditions receive the same change inventory for the crop being described, and
CFS scores only claims about \emph{that} crop; retrieved history concerns
\emph{other} crops, so it can move wording but rarely a claim. Measured
directly: of \EFourNCrops{} generations, \EFourTextDiffer{}
(\EFourTextDifferPct\%) differ in text between the three conditions and
only \EFourClaimsDiffer{} (\EFourClaimsDifferPct\%) differ in extracted
claims. A crop-level ladder therefore cannot decide the question it appears to
ask, and reporting the three as an ablation would be presenting noise as a
result. Section~\ref{ssec:e7} asks it at the unit where the answer can differ.

\textbf{Generation adds no measurable factuality over the deterministic
template.} The template states its claims from the same inventory without a
language model and is not separated from the generated conditions on CFS. What
generation adds is prose a subscriber can read, which is a different claim
requiring different evidence, and we do not measure readability here.

\subsection{Pairing Determines Report Factuality (E5)}

E4 varies grounding with pairing fixed; E5 does the reverse. Grounding is
pinned at \texttt{llm\_graphrag} and only the pairing is swapped, so whatever
moves is attributable to pairing. Unlike the graph term, the swapped variable
demonstrably reaches the evidence: every crop receives a different change
inventory.

% generated by paper/make_tables.py -- do not edit
% CLIP + geometry heuristic: results/levir_cc_caption_heuristic_pairing (n=1929, checkpoint=None)
% learned head (ours): results/levir_cc_caption (n=1929, checkpoint=data/checkpoints/pairing_head.pt)
\begin{table}[t]
\centering
\caption{Report factuality under two pairing methods (E5).}
\label{tab:e5_pairing}
\setlength{\tabcolsep}{4pt}
\begin{tabular}{lccccc}
\toprule
Pairing & CFS-P & CFS-R & CFS-F1 & Hal & ChgAcc \\
\midrule
CLIP + geometry heuristic & 34.01 & 36.79 & 35.35 & 65.99 & 65.01 \\
\textbf{learned head (ours)} & \textbf{61.72} & \textbf{39.70} & \textbf{48.32} & \textbf{38.28} & \textbf{76.52} \\
\bottomrule
\end{tabular}
\vspace{1pt}
\parbox{\columnwidth}{\scriptsize\textit{Note:} Grounding is fixed at \texttt{llm\_graphrag}; crops, style examples and the language model are identical across rows. The heuristic has no verifier, so this gap combines matching and suppression; Section~\ref{ssec:ablation} decomposes them for detection.}
\end{table}


Report factuality moves by \EFiveGain{} CFS-F1. Two of every three claims the
heuristic-fed report makes are unsupported (\EFiveHeuristicHal\%), against one
in four (\EFiveLearnedHal\%) for the learned head; ChgAcc rises from
\EFiveHeuristicChgAcc{} to \EFiveLearnedChgAcc. This is the measurement that
connects detection quality to the product: E1 and E2 show the head detects
better, and E5 shows the improvement survives into what the subscriber is
shown. A paper with E1 alone would be asserting that link rather than measuring
it.

\subsection{Neighbourhood-Level Grounding (E7)}
\label{ssec:e7}

At crop level the conditions are nested and the graph term vanishes. At
neighbourhood level they stop being a ladder and become three
\emph{representations} of the same observations --- full concatenation, top-$k$
retrieval, graph aggregate --- and the comparison becomes decidable. One report
is generated per neighbourhood over every crop of the tile the split provides.

% generated by paper/make_tables.py -- do not edit
% source: results/levir_cc_scene (checkpoint=data/checkpoints/pairing_head.pt)
\begin{table}[t]
\centering
\caption{Neighborhood-level grounding (E7): one report per neighborhood, over every crop of a tile the split provides. The conditions are three representations of the same observations, not an additive ladder. $n=218$, 1--16 crops each (mean 8.8; 129 scenes with more than one crop), mean 37.7 ground-truth instances per scene.}
\label{tab:e7_scene}
\setlength{\tabcolsep}{4pt}
\begin{tabular}{lccc}
\toprule
Representation & CFS-F1 & Hal & scene CountMAE \\
\midrule
template (no LLM) & 40.99 & 56.46 & 26.15 \\
full concatenation & 21.02 & 76.58 & 458.46 \\
top-$k$ retrieval & 18.54 & 78.57 & 167.87 \\
graph aggregate & 27.55 & 72.28 & 18.50 \\
\bottomrule
\end{tabular}
\end{table}


\textbf{Counting is where representation decides the outcome, and the margins
are not subtle.} Concatenating every crop's inventory into one prompt produces
a scene count error of \ESevenStructCountMAE{} instances; top-$k$ retrieval,
which drops crops by construction, reaches \ESevenFlatCountMAE; graph
aggregation reaches \ESevenGraphCountMAE. The failure mode of concatenation is
worth naming: handed everything, the language model does not aggregate it, and
the error grows with the number of observations rather than shrinking. Top-$k$
retrieval is better only because it truncates the prompt, and it cannot recover
a count from crops it never retrieved. Aggregating before generation is what
survives.

\textbf{Graph aggregation is the only generated condition that beats the
deterministic template on any axis}, and it does so on the axis this
application cares about: \ESevenGraphCountMAE{} against the template's
\ESevenTemplateCountMAE. On claim-level factuality it leads the other generated
conditions (\ESevenGraphFOne{} against \ESevenStructFOne{} and
\ESevenFlatFOne) but still trails the template (\ESevenTemplateFOne). We
therefore claim that graph aggregation preserves counts better than either
retrieval or concatenation, and that it is what makes generated
neighbourhood-level reports numerically usable --- not that generation improves
factuality overall.

\textbf{Scale is what made this measurable.} An earlier run over
\PilotESevenNScenes{} scenes put graph aggregation \emph{below} plain
concatenation on CFS-F1 (\PilotESevenGraphFOne{} against
\PilotESevenStructFOne). Over \ESevenNScenes{} scenes the ordering reverses and
the counting gap widens by more than an order of magnitude. The two results are
consistent rather than contradictory: with few observations any representation
serves, and the cost of not aggregating appears only once there are enough
observations to lose track of.

\subsection{Deployment Cost (E6)}

\begin{table}[t]
\centering
\caption{Per-tile processing cost on a single NVIDIA A100 80GB. The learned pairing head adds 19781 parameters and 0.70~ms/tile over the hand-tuned heuristic it replaces --- 0.01\% of the pipeline. Segmentation and language generation dominate the budget.}
\label{tab:efficiency}
\begin{tabular}{lrr}
\toprule
Stage & Latency (ms/tile) & Peak GPU (MB) \\
\midrule
SAM3 detection + CLIP (2 frames) & 6099.45 & -- \\
pairing: heuristic (production) & 8.98 & 0 \\
pairing: learned head (ours) & 9.69 & 10 \\
knowledge graph: index + retrieve & 141.56 & 8 \\
report LLM (server:LGAI-EXAONE/EXAONE-4.0-32B@http://127.0.0.1:8005/v1, batched) & 438.29 & 8 \\
\bottomrule
\end{tabular}
\end{table}


Measured over 512 pairs of the same cache the report experiments use.
Segmentation is \ESixDetectShare\% of the per-tile budget. The learned head
costs \ESixHeadMs~ms against \ESixHeuristicMs~ms for the heuristic it replaces
--- \ESixHeadDelta~ms more, \ESixPipelineShare\% of the pipeline. We quote the
difference rather than characterising it, because the two rows it is derived
from are in the table above and a reader can check the subtraction.

The comparison is only meaningful within one cache. The same measurement over
the 8-scene pilot cache gives 1.45~ms, because those crops carry a different
number of detections and both pairing strategies scale with that count; two
runs over the full cache agree to 0.08~ms. A per-tile latency for this stage is
therefore a property of the tile population, not a constant of the method, and
we report it against the population the rest of the paper uses.
